\begin{figure}[t]
\centering
\begin{adjustbox}{max width=\textwidth,center}
\begin{tikzpicture}[
  node distance=7mm and 7mm,
  sector/.style={draw,rounded corners,fill=gray!7,align=center,
    minimum height=12mm,text width=0.19\textwidth,font=\small},
  equation/.style={draw,very thick,rounded corners,align=center,
    minimum height=13mm,text width=0.31\textwidth,font=\small},
  endpoint/.style={draw,double,rounded corners,align=center,
    minimum height=13mm,text width=0.27\textwidth,font=\small},
  note/.style={draw,dashed,rounded corners,align=left,
    text width=0.27\textwidth,font=\scriptsize},
  arrow/.style={-{Latex[length=2.2mm]},thick}
]
\node[sector] (h1) {conifold hyper 1\\$u_1=h_1\widetilde h_1$};
\node[sector,right=of h1] (h2) {conifold hyper 2\\$u_2=h_2\widetilde h_2$};
\node[sector,right=of h2] (e2) {$E_2$ fixed point\\$\beta=\mu^{\CC}_{SU(2)}$};

\node[equation,below=11mm of h2] (mm)
  {$Q^{\mathrm{eff}}_{X_9}=\begin{psmallmatrix}1&0&1\\0&1&1\end{psmallmatrix}$\\[1mm]
   $u_1+\beta=0$\,,\quad $u_2+\beta=0$};
\node[endpoint,below=9mm of mm] (smooth)
  {primitive mixed trajectory\\
   $(u_1,u_2,\beta)=t(-1,-1,1)$};
\node[note,right=11mm of mm] (vectors)
  {$V_1,V_2$: flavor vectors Higgsed by the rank-two gauging.\\[1mm]
   $E$: neutral $E_2$ Coulomb vector lost when the exceptional divisor is
   smoothed; it is not a third gauging.};

\draw[arrow] (h1.south) -- ++(0,-5mm) -| ([xshift=-6mm]mm.north);
\draw[arrow] (h2.south) -- (mm.north);
\draw[arrow] (e2.south) -- ++(0,-5mm) -| ([xshift=6mm]mm.north);
\draw[arrow] (mm.south) -- node[right,font=\scriptsize] {$\ker Q_{X_9}^{\mathrm{eff}}$} (smooth.north);
\draw[arrow,dashed] (mm.east) -- (vectors.west);
\end{tikzpicture}
\end{adjustbox}
\caption{Cooperative Higgsing in $X_9$. No nonzero trajectory is supported on
the two conifold sectors or on the $E_2$ sector alone. The unique global
direction activates all three, while the protected vector-multiplet count
separates the two Higgsed flavor vectors from the neutral SCFT Coulomb vector.}
\label{fig:x9-cooperation}
\end{figure}
